\documentclass[11pt]{article}

\usepackage[margin=1in]{geometry}
\usepackage{amsmath,amssymb,amsthm}
\usepackage{mathtools}
\usepackage{enumitem}
\usepackage{hyperref}

\newtheorem{theorem}{Theorem}[section]
\newtheorem{proposition}[theorem]{Proposition}
\newtheorem{lemma}[theorem]{Lemma}
\newtheorem{corollary}[theorem]{Corollary}
\theoremstyle{definition}
\newtheorem{remark}[theorem]{Remark}
\newtheorem{program}[theorem]{Program}
\newtheorem{target}[theorem]{Target Lemma}

\newcommand{\Tr}{\operatorname{Tr}}
\newcommand{\1}{\mathbf{1}}
\newcommand{\HS}{\mathrm{HS}}
\newcommand{\R}{\mathbb{R}}

\title{New Strategies Toward the Odd-Cycle Goodman Bound in Graphons\\
\large A companion note: one reformulation, two uniform criteria, and two programs}
\author{Working note}
\date{}

\begin{document}
\maketitle

\begin{abstract}
This note accompanies the audited working note proving the odd-cycle lower bound
$t(C_{2k+1},W)\ge p^{2k+1}-p(1-p)^{2k}$ for $C_3,\dots,C_{13}$ and under various
structural hypotheses. We pursue strategies aimed at \emph{all} odd cycles
simultaneously, and aimed at statements of independent interest rather than
merely re-deriving the regular case.

We contribute three rigorous results and set up two programs.
\textbf{(R1)} A generating-function reformulation: the conjecture for all odd $m$
is equivalent to the coefficientwise nonnegativity of a single resolvent
difference $\Phi_W(w)$, with the Perron term separating off as a coefficientwise
nonnegative summand; the residual positivity is \emph{tight at bipartite},
which both explains why the Perron/non-principal coupling is unavoidable and
recasts the problem as a Markov-moment (Hankel) domination.
\textbf{(R2)} A uniform spectral criterion phrased intrinsically on $T_W$:
if $T_W\succeq-(1-p)I$ (equivalently $\lambda_{\min}(T_W)\ge-q$) then the
conjecture holds for every odd $m$; the proof is four lines and does not invoke
the hub-coupling lemma. For irregular graphons this is a genuinely different
sufficient condition than the positive-spectrum-safe criterion, with which it
coincides in the regular case.
\textbf{(R3)} For $p\ge 2/3$ the trace-square budget forces at most one
eigenvalue of $T_W$ below $-q$, reducing the entire regime to a two-parameter
scalar family and a single inequality $(\star)$; this isolates the precise
missing ingredient as a lower bound on the Perron surplus in terms of the deep
mode depth.
Finally we set up \textbf{(P1)} a reflection-positivity / transfer-operator
symmetrization program and \textbf{(P2)} a Markov--Krein resolvent-domination
program, the two routes that can in principle be uniform in $m$.

We are explicit about status throughout: R1--R3 are proved; the obstruction
discussion and the two programs are clearly marked as conjectural.
\end{abstract}

\section{Setup and notation}\label{sec:setup}

Throughout $W:[0,1]^2\to[0,1]$ is a graphon, $p=t(K_2,W)$ is its edge density,
$q=1-p$, and $T=T_W$ is the associated self-adjoint Hilbert--Schmidt operator
$T_Wf(x)=\int_0^1 W(x,y)f(y)\,dy$. We work in the nontrivial regime $p>1/2$,
i.e.\ $0\le q<1/2$, since $g_m(p):=p^m-p\,q^{m-1}\le 0$ for $p\le 1/2$. Set
$m=2k+1$.

We use the spectral facts established in the companion note; for completeness we
restate them.
\begin{enumerate}[label=(S\arabic*)]
\item For $m\ge 2$, $t(C_m,W)=\Tr(T_W^m)=\sum_i\lambda_i^m$, where
$\lambda_0\ge\lambda_1\ge\cdots$ are the eigenvalues of $T_W$ with multiplicity.
\item (Perron / Rayleigh.) The top eigenvalue $\lambda_0=:\ell$ is nonnegative
and $\ell\ge\langle T_W\1,\1\rangle=p$.
\item (Second-moment budget.) $\sum_i\lambda_i^2=\Tr(T_W^2)=\|W\|_2^2\le p$, hence
$\ell\le\sqrt p$ and, using $\ell\ge p$,
\[
\sum_{i\ge1}\lambda_i^2=\|W\|_2^2-\ell^2\le p-\ell^2\le p-p^2=pq.
\]
\item (Spectral support; companion Lemma~8.7.) Every non-principal eigenvalue of
$T_W$ lies in $[-\tfrac12,\tfrac12]$; in particular $\lambda_{\min}(T_W)\ge-\tfrac12$.
\item (Triangle density; Razborov--Reiher, used only in
\S\ref{sec:reduction}.) For $\tfrac12\le p\le\tfrac23$, with $c\in[0,\tfrac13]$
defined by $p=\tfrac12+c-\tfrac32c^2$, one has $t(K_3,W)\ge\Theta(p):=\tfrac32c(1-c)^2$.
\end{enumerate}

We also use the block decomposition of the complement $U=1-W$, $T_U=J-T_W$ where
$Jf=\langle f,\1\rangle\1$, with degree fluctuation $g=T_U\1-q\1\in\1^\perp$ and
mean-zero compression $A=P_{\1^\perp}T_UP_{\1^\perp}$; thus
\[
T_U=\begin{pmatrix} q & g^*\\ g & A\end{pmatrix},
\qquad
T_W=\begin{pmatrix} p & -g^*\\ -g & -A\end{pmatrix}
\quad\text{on }\R\1\oplus\1^\perp .
\]
The graphon is regular iff $g=0$.

\section{A generating-function reformulation}\label{sec:genfun}

We package all odd cycles into a single analytic object. For $\lambda\in[-1,1]$,
\[
\sum_{k\ge0}\lambda^{2k+1}w^k=\frac{\lambda}{1-\lambda^2w}\qquad(|w|<\lambda^{-2}),
\]
and the target sequence has the closed form
\[
\sum_{k\ge0}g_{2k+1}(p)\,w^k
=\sum_{k\ge0}\bigl(p^{2k+1}-p\,q^{2k}\bigr)w^k
=\frac{p}{1-p^2w}-\frac{p}{1-q^2w}.
\]

\begin{theorem}[Resolvent reformulation]\label{thm:reform}
Define, for $|w|<4$,
\[
\Phi_W(w)\;:=\;\sum_i\frac{\lambda_i}{1-\lambda_i^2w}
-\frac{p}{1-p^2w}+\frac{p}{1-q^2w}.
\]
Then $\Phi_W$ is analytic near $0$ and
\[
[w^k]\,\Phi_W \;=\; t(C_{2k+1},W)-g_{2k+1}(p)\qquad(k\ge0).
\]
Consequently Conjecture~1.1 holds for $W$ for all odd $m\ge3$ if and only if
$\Phi_W$ has nonnegative Taylor coefficients at $w=0$.
\end{theorem}

\begin{proof}
By (S1)--(S4) the eigenvalues satisfy $|\lambda_i|\le 1$ and
$\sum_i\lambda_i^2\le p<\infty$, so $\sum_i\frac{\lambda_i}{1-\lambda_i^2w}$
converges locally uniformly for $|w|<1$ (indeed $|w|<4$ off the principal mode),
and $\Phi_W$ is analytic near $0$. Expanding each summand,
\[
[w^k]\sum_i\frac{\lambda_i}{1-\lambda_i^2w}=\sum_i\lambda_i^{2k+1}=t(C_{2k+1},W),
\]
using (S1) for $2k+1\ge3$ (and $k=0$ gives $\sum_i\lambda_i=\Tr T_W=p=t(K_2,W)$,
matching $g_1=p$). Subtracting the closed forms above gives the claimed
coefficient identity, and the equivalence is immediate.
\end{proof}

\paragraph{The Perron split.}
Separate the principal mode $\ell=\lambda_0\ge p$:
\begin{equation}\label{eq:split}
\Phi_W(w)=\underbrace{\Bigl(\frac{\ell}{1-\ell^2w}-\frac{p}{1-p^2w}\Bigr)}_{=:P(w)}
\;+\;\underbrace{\sum_{i\ge1}\frac{\lambda_i}{1-\lambda_i^2w}+\frac{p}{1-q^2w}}_{=:R(w)} .
\end{equation}
Since $\ell\ge p$, $[w^k]P=\ell^{2k+1}-p^{2k+1}\ge0$: \emph{the Perron term is
coefficientwise nonnegative}. It is tempting to conclude that it suffices to
prove $R$ coefficientwise nonnegative, i.e.\
$\sum_{i\ge1}\lambda_i^{2k+1}\ge -p\,q^{2k}$. We record why this lossy split
\emph{cannot} be the whole story, as it is instructive.

\begin{remark}[Tightness at bipartite; the split is lossy]\label{rem:tight}
At a balanced complete bipartite graphon ($p=\tfrac12$, $\ell=\tfrac12$, one
non-principal eigenvalue $-\tfrac12$, rest $0$) one has, for every $k$,
\[
\sum_{i\ge1}\lambda_i^{2k+1}=-(\tfrac12)^{2k+1}=-p\,q^{2k},
\]
so $R$ has \emph{zero} $w^k$-coefficient: the residual inequality is exactly
tight, with the Perron surplus $[w^k]P=0$. Perturbing $p$ slightly above
$\tfrac12$ while keeping a deep mode near $-\tfrac12$ would make
$\sum_{i\ge1}\lambda_i^{2k+1}\approx-(\tfrac12)^{2k+1}<-pq^{2k}$, violating the
residual inequality; such spectra are forbidden by $0\le W\le1$ (they are not
graphon-realizable), but only the Perron surplus $[w^k]P>0$ in the genuine
graphon redeems them. Thus \eqref{eq:split} cannot be decoupled: the Perron and
non-principal parts must be used \emph{together}. This is the
generating-function incarnation of the ``aggregate, not modewise'' obstruction
identified in companion \S9.6.
\end{remark}

\begin{remark}[Even cycles are free]\label{rem:even}
Replacing $\frac{\lambda}{1-\lambda^2w}$ by $\frac{\lambda^2 w}{1-\lambda^2 w}$
gives the even part, whose coefficients are $\sum_i\lambda_i^{2k}=t(C_{2k},W)\ge
p^{2k}\ge g_{2k}(p)$ by bipartite Sidorenko. Hence \emph{all} the content of the
even-and-odd statement is carried by the odd part $\Phi_W$; the reformulation
loses nothing.
\end{remark}

The value of Theorem~\ref{thm:reform} is methodological: it turns the problem
into coefficientwise (Hankel/total-positivity) domination of one Stieltjes-type
function by a two-pole rational function, which is the natural input for
Markov--Krein theory (\S\ref{sec:markov}).

\section{A uniform criterion: \texorpdfstring{$T_W\succeq-(1-p)I$}{TW >= -(1-p)I}}
\label{sec:negsafe}

We give a clean sufficient condition, phrased intrinsically on the least
eigenvalue of $T_W$, valid for every odd cycle at once.

\begin{theorem}[Negative-spectrum-safe criterion]\label{thm:negsafe}
Let $p>\tfrac12$. If
\[
\lambda_{\min}(T_W)\ \ge\ -q\qquad\text{(equivalently }T_W+qI\succeq0\text{)},
\]
then for every odd $m\ge3$,
\[
t(C_m,W)\ \ge\ p^m-p\,q^{m-1}=g_m(p).
\]
\end{theorem}

\begin{proof}
By (S1), splitting the spectrum into the Perron mode $\ell=\lambda_0$, the
positive non-principal eigenvalues, and the negative ones (write
$\mu_j=-\lambda_j>0$ for $\lambda_j<0$),
\[
t(C_m,W)=\ell^m+\!\!\sum_{i\ge1,\,\lambda_i>0}\!\!\lambda_i^m-\sum_j\mu_j^m
\ \ge\ \ell^m-\sum_j\mu_j^m,
\]
since $m$ is odd and the positive terms are nonnegative. By hypothesis every
$\mu_j\le q$, so $\mu_j^m\le q^{m-2}\mu_j^2$ (as $m-2\ge1$), whence
\[
\sum_j\mu_j^m\ \le\ q^{m-2}\sum_j\mu_j^2\ \le\ q^{m-2}\sum_{i\ge1}\lambda_i^2
\ \le\ q^{m-2}\,pq\ =\ p\,q^{m-1},
\]
using the budget (S3). Therefore $t(C_m,W)\ge\ell^m-pq^{m-1}\ge p^m-pq^{m-1}$,
the last step from $\ell\ge p$ (S2).
\end{proof}

\begin{remark}[Comparison with known criteria]\label{rem:compare}
For a \emph{regular} graphon $g=0$, so $T_W$ block-diagonalizes and on $\1^\perp$
one has $T_W=-A$; then $\lambda_{\min}(T_W)\ge-q\iff\lambda_{\max}(A)\le q$, the
positive-spectrum-safe criterion. Thus Theorem~\ref{thm:negsafe} contains the
regular case, recovering companion Corollary~9.6. For an \emph{irregular}
graphon ($g\ne0$) the two criteria differ: $\lambda_{\min}(T_W)$ feels the
coupling $g$ between the Perron mode and $\1^\perp$ through the off-diagonal
block in $T_W$, whereas $\lambda_{\max}(A)$ is a property of the compressed
complement alone. Neither implies the other in general; they are complementary
sufficient conditions. The advantage of Theorem~\ref{thm:negsafe} is that it is
a single positive-semidefiniteness statement about $T_W$ with a four-line proof.
\end{remark}

\begin{remark}[Sharpness and scope]\label{rem:scope}
The constant $-q$ cannot be relaxed to $-\tfrac12$: at $p$ slightly above
$\tfrac12$ a near-bipartite graphon has $\lambda_{\min}(T_W)$ near $-\tfrac12<-q$
and is exactly the hard case. Thus the criterion covers all graphons whose
spectrum is ``not too negative'' (a strict superset of regular graphons,
uniform in $m$), and fails precisely on the near-bipartite family that the
companion note also identifies as the genuine obstruction near $p=\tfrac12$.
\end{remark}

\section{The regime \texorpdfstring{$p\ge 2/3$}{p>=2/3}: reduction to a scalar
inequality}\label{sec:reduction}

We now show that for $p\ge 2/3$ the whole problem reduces to a two-parameter
scalar family, and we identify the single missing ingredient.

\begin{lemma}[At most one deep mode]\label{lem:onedeep}
Let $p\ge 2/3$, so $q\le\tfrac13$. Then $T_W$ has at most one eigenvalue strictly
below $-q$.
\end{lemma}

\begin{proof}
An eigenvalue below $-q$ has square exceeding $q^2$. Two of them would contribute
more than $2q^2$ to $\sum_{i\ge1}\lambda_i^2$, but the budget (S3) gives
$\sum_{i\ge1}\lambda_i^2\le pq$. Since $p\ge2/3$ means $2q\le p$, i.e.\
$2q^2\le pq$, two deep modes are impossible.
\end{proof}

\begin{proposition}[One-deep-mode reduction]\label{prop:onedeep}
Let $p\ge 2/3$. Write $\ell=\lambda_0\ge p$ for the Perron eigenvalue and, if
$T_W$ has an eigenvalue below $-q$, let $-\eta$ be that (unique, by
Lemma~\ref{lem:onedeep}) eigenvalue, with $\eta\in(q,\tfrac12]$ by (S4);
otherwise the hypothesis of Theorem~\ref{thm:negsafe} holds and we are done.
If, for every odd $m\ge3$,
\begin{equation}\label{eq:star}
\bigl(\ell^m-p^m\bigr)+q^{m-2}\bigl(\ell^2-p^2\bigr)
\ \ge\ \eta^2\bigl(\eta^{m-2}-q^{m-2}\bigr),
\tag{$\star$}
\end{equation}
then $t(C_m,W)\ge g_m(p)$ for all odd $m\ge3$.
\end{proposition}

\begin{proof}
As in Theorem~\ref{thm:negsafe}, drop the positive non-principal eigenvalues:
\[
t(C_m,W)\ \ge\ \ell^m-\eta^m-\sum_{\text{other neg.}}\mu_j^m.
\]
The negative eigenvalues other than $-\eta$ have magnitude $\le q$ and, by the
budget (S3), total square at most $\;p-\ell^2-\eta^2$. Hence
$\sum_{\text{other}}\mu_j^m\le q^{m-2}\bigl(p-\ell^2-\eta^2\bigr)$, so
\[
t(C_m,W)\ \ge\ \ell^m-\eta^m-q^{m-2}\bigl(p-\ell^2-\eta^2\bigr).
\]
Using $p\,q^{m-1}=q^{m-2}(p-p^2)$, the desired bound $t(C_m,W)\ge p^m-pq^{m-1}$
rearranges exactly to \eqref{eq:star}.
\end{proof}

\begin{remark}[What \eqref{eq:star} reduces to]\label{rem:starsmallm}
When $\eta\le q$ the right side of \eqref{eq:star} is $\le0$ and the inequality
holds trivially: this recovers Theorem~\ref{thm:negsafe}. Since $p\ge2/3>1/2\ge
\eta$, we have $\ell\ge p>\eta$, so $\ell^m-\eta^m$ grows and the left side
dominates for large $m$; the binding constraints are the smallest odd $m$. Thus
\eqref{eq:star} for all odd $m$ is, in practice, a check at $m=3,5$. The full
$p\ge2/3$ region is therefore reduced to a two-dimensional scalar question in
$(\ell,\eta)$ subject to the realizability constraints $\ell\ge p$,
$\eta\le\tfrac12$, and $\ell^2+\eta^2\le p$.
\end{remark}

\begin{remark}[The residual obstruction, made precise]\label{rem:obstruction}
Proposition~\ref{prop:onedeep} is \emph{not} unconditional, and the gap is
exactly the coupling obstruction. The worst case for \eqref{eq:star} is
$\ell=p$ (no Perron surplus) with $\eta$ as large as the budget allows
($\eta=\sqrt{pq}$), where the left side vanishes but the right side is strictly
positive. However, $\ell=p$ forces $g=0$ (regularity), and a regular graphon has
\emph{all} non-principal eigenvalues bounded by $q$ in modulus, so it has no deep
mode at all: the configuration that defeats \eqref{eq:star} is not
graphon-realizable. Realizing a deep mode $-\eta$ with $\eta>q$ requires
irregularity ($g\ne0$), which in turn raises $\ell$ above $p$ through the
off-diagonal block of $T_W$. The missing quantitative content is precisely a
graphon lower bound on the Perron surplus in terms of the deep-mode depth:
\end{remark}

\begin{target}[Perron surplus vs.\ deep mode]\label{tgt:surplus}
There is a function $\sigma(\eta;p)$ with $\sigma(\eta;p)\ge 0$ and
$\sigma(q;p)=0$ such that every graphon $W$ with $p\ge 2/3$ and a deep mode
$-\eta$ of $T_W$ ($\eta\in(q,\tfrac12]$) satisfies
\[
\bigl(\ell^2-p^2\bigr)\;+\!\!\sum_{i\ge1,\,\lambda_i>0}\!\!\lambda_i^2
\ \ge\ \sigma(\eta;p),
\]
with $\sigma$ large enough that \eqref{eq:star} holds for all odd $m\ge3$.
\end{target}

Target Lemma~\ref{tgt:surplus} would close $p\ge 2/3$ uniformly in $m$ via
Proposition~\ref{prop:onedeep}. It is a statement purely about the second-order
geometry of $T_W$ (no high cycle powers), and it is exactly the
``aggregate-not-modewise'' content of companion \S9.6: when the deep mode
decouples from $g$ (so $\ell$ gets no help from it), the depth $\eta$ must itself
be constrained by $0\le W\le1$ through the orthogonality of its eigenvector to
the constants. Quantifying that constraint is, in our view, the cleanest
finite-dimensional open problem in this circle of ideas.

\section{Program P1: reflection positivity / transfer-operator symmetrization}
\label{sec:rp}

The quantity $t(C_m,W)=\Tr(T_W^m)$ is the periodic partition function of the
one-dimensional transfer operator $T_W$ on $m$ sites, and the cyclic/reflection
symmetry of $C_m$ is \emph{the same for every $m$}. This is the one symmetry that
is intrinsically uniform in $m$, which is why reflection positivity (RP) is the
natural tool for a uniform structural theorem. We set up what one would need.

\paragraph{The cut and the gluing form.}
Fix $m=2k+1$. Cut $C_m$ at one vertex into a path $P_{m-1}$ with two marked
endpoints, and consider the bilinear pairing on $L^2([0,1])$,
\[
\langle f,g\rangle_{T^{\,m-1}}:=\int_{[0,1]^2}(T_W^{\,m-1})(x,y)\,f(x)g(y)\,dx\,dy
=\bigl\langle f,\,T_W^{\,m-1}g\bigr\rangle .
\]
Because $m-1$ is even, $T_W^{\,m-1}\succeq0$, so this pairing is positive
semidefinite: it is a Gram form. Taking $f=g=\1$ recovers
$t(P_{m-1},W)=\langle\1,T_W^{m-1}\1\rangle$, and gluing the two endpoints back
(integrating against $W$) recovers $t(C_m,W)$. The RP content is that the
even-length ``half-cycle'' folds positively.

\paragraph{The symmetrization conjecture.}
Let $\mathcal{W}(p,\tau)=\{W:t(K_2,W)=p,\ \Tr(T_W^2)=\tau\}$. The program is to
prove a one-step symmetrization inequality: there is an involution-based
operation $W\mapsto W^{\mathrm{sym}}$ (folding $W$ across the spectral
reflection $\lambda\mapsto$ its RP-conjugate) that preserves $(p,\tau)$ and does
not increase $\Tr(T_W^m)$ for any odd $m$, with fixed points exactly the
reflection-symmetric (near-bipartite) graphons.

\begin{program}[RP symmetrization]\label{prog:rp}
Establish that on $\mathcal{W}(p,\tau)$ the minimizer of $t(C_m,\cdot)$ may be
taken reflection-symmetric for every odd $m$ simultaneously; equivalently, that
the negative part of the spectral measure of $T_W$ may be concentrated at a
single reflected atom without increasing any odd trace power. Then close the
reduced (symmetric) family by the low-dimensional inequality of
\S\ref{sec:reduction}.
\end{program}

\paragraph{Why it can be uniform, and where it is delicate.}
The folding operation and the PSD form above are defined identically for all
$m$, so a successful symmetrization step is automatically uniform; this is the
structural analogue of ``the minimizer near $p=\tfrac12$ is near-bipartite,''
which the companion note observes numerically (\S10.5) but does not prove. The
delicacy is genuine: odd cycles are \emph{not} bipartite, so the naive single
cut does not fold cleanly, and RP for odd cycles typically yields PSD
\emph{constraints} (Cauchy--Schwarz inequalities between path overlaps) that
still require assembly into a monotonicity statement. The first concrete
sub-target is the odd-cut inequality
\[
t(C_{2k+1},W)^2\ \le\ t(C_{2k},W)\cdot t(C_{2k+2},W)\quad(?),
\]
a log-convexity of the cycle sequence that would follow from a single RP fold and
would, with the even-cycle Sidorenko tower (Remark~\ref{rem:even}), constrain the
odd cycles from the (known) even ones.

\section{Program P2: Markov--Krein resolvent domination}\label{sec:markov}

Theorem~\ref{thm:reform} recasts the conjecture as coefficientwise domination of
the Stieltjes-type function $\sum_i\frac{\lambda_i}{1-\lambda_i^2w}$ by the
two-pole function $\frac{p}{1-p^2w}-\frac{p}{1-q^2w}$. Coefficientwise (rather
than pointwise) domination is a moment problem, and Markov--Krein / canonical
moment theory is built to reduce such problems to finitely many atoms.

\paragraph{The moment body.}
Regard the spectrum as a positive measure $\rho=\sum_i\delta_{\lambda_i}$ on
$[-\tfrac12,\tfrac12]$ (after separating the Perron atom $\ell$). The genuine
graphon constraints translate to moment constraints on $\rho$:
\[
\textstyle\int 1\,d\rho_{\ge1}=:M_0,\quad \int\lambda^2\,d\rho_{\ge1}=M_2\le pq,
\quad \int\lambda^3\,d\rho_{\ge1}=t(K_3,W)-\ell^3\ge\Theta(p)-\ell^3,
\]
together with the even Sidorenko bounds
$\int\lambda^{2j}\,d\rho_{\ge1}=t(C_{2j},W)-\ell^{2j}\ge p^{2j}-\ell^{2j}$.
Minimizing $\int\lambda^{m}\,d\rho_{\ge1}$ subject to a \emph{fixed finite} set
of these constraints, the Markov--Krein theorem returns a canonical
representation with atom count controlled by the number of binding constraints,
\emph{not} by $m$. That is the rigorous form of ``reduce to a low-dimensional
family,'' and it is uniform in $m$ by construction.

\begin{program}[Matrix Markov--Krein on $(g,A)$]\label{prog:mk}
Carry out the canonical-moment reduction, but on the \emph{joint} moments
$s_j=\langle g,A^jg\rangle$ of the pair $(g,A)$ rather than on the marginal
spectral moments of $\rho$. The marginal-moment problem is provably defeatable
(companion \S9.6: a decoupled deep mode satisfies all marginal-moment
constraints yet the conjecture survives via the coupling), so the correct object
is the operator/matrix moment problem in which $g$ and $A$ appear jointly. The
deliverable is: among graphon-admissible $(g,A)$ with bounded $s_0,s_1,s_2$, the
minimizer of $\Tr(T_W^m)$ is a low-rank (arrowhead) configuration for all odd $m$.
\end{program}

\paragraph{Status.}
The reduction to finitely many atoms is rigorous \emph{given} which constraints
bind; the open content is (i) identifying the binding set and (ii) lifting from
marginal to joint moments so that the coupling $g$ is captured. We emphasize the
honest ceiling established in the companion note: \emph{any} argument using a
\emph{fixed} finite set of marginal moments has a window in $p$ shrinking like
$m^{-2/(r-2)}$ for an $r$-th moment input, hence cannot be uniform; Program P2
escapes this only through the joint-moment lift, which is where the real
difficulty migrates.

\section{Program P3: a uniform telescoping recurrence}\label{sec:telescope}

The solved structural families (rank-one \S9.4, equal cliques \S9.5, two-value
multipartite \S10.3 of the companion note) all close \emph{for all odd $m$ at
once} by the same device: a two-step recurrence inequality with a nonnegative
increment, seeded by a base case. We abstract this and identify exactly the
general obstruction it meets. The point is that the recurrence parameter $a=q$ is
canonical, because it annihilates the target's correction term.

\begin{lemma}[Canonical telescoping identity]\label{lem:tele}
Let $D_m:=t(C_m,W)-g_m(p)$. For every $m\ge3$,
\begin{equation}\label{eq:tele}
D_{m+2}-q^2D_m=\sum_i\lambda_i^m\bigl(\lambda_i^2-q^2\bigr)-p^m\bigl(p^2-q^2\bigr).
\end{equation}
Moreover the Perron term covers the right-hand subtrahend:
$\ell^m(\ell^2-q^2)\ge p^m(p^2-q^2)$ (since $\ell\ge p>q$), so
\begin{equation}\label{eq:tele2}
D_{m+2}-q^2D_m\ \ge\ \sum_{i\ge1}\lambda_i^m\bigl(\lambda_i^2-q^2\bigr).
\end{equation}
\end{lemma}

\begin{proof}
Write $D_m=\sum_i\lambda_i^m-p^m+pq^{m-1}$. Then
$D_{m+2}-q^2D_m=\sum_i\lambda_i^m(\lambda_i^2-q^2)-(p^{m+2}-q^2p^m)
+(pq^{m+1}-q^2pq^{m-1})$. The last bracket is $pq^{m+1}-pq^{m+1}=0$, and
$p^{m+2}-q^2p^m=p^m(p^2-q^2)$, giving \eqref{eq:tele}. For \eqref{eq:tele2}, the
$i=0$ term is $\ell^m(\ell^2-q^2)$ with $\ell\ge p$ and $\ell^2-q^2\ge p^2-q^2>0$
(as $p>1/2>q$ implies $p>q$), so $\ell^m(\ell^2-q^2)\ge p^m(p^2-q^2)$.
\end{proof}

\begin{proposition}[Telescoping criterion]\label{prop:tele}
Suppose $D_3\ge0$ and, for every odd $m\ge3$,
\begin{equation}\label{eq:teleincr}
\sum_{i\ge1}\lambda_i^m\bigl(\lambda_i^2-q^2\bigr)\ \ge\ 0 .
\end{equation}
Then $D_m\ge0$ for all odd $m\ge3$, i.e.\ the conjecture holds for $W$.
\end{proposition}

\begin{proof}
By \eqref{eq:tele2} and \eqref{eq:teleincr}, $D_{m+2}\ge q^2D_m\ge0$ whenever
$D_m\ge0$; induction from $D_3\ge0$ gives all odd $m$.
\end{proof}

\paragraph{What \eqref{eq:teleincr} sees.}
The summand $\lambda^m(\lambda^2-q^2)$ is nonnegative exactly when $\lambda\ge q$
or $-q\le\lambda\le0$, and negative for $\lambda\in(0,q)$ (small positive modes)
and for $\lambda<-q$ (deep negative modes). Thus the telescoping increment is
governed by the \emph{same} two enemies that appear everywhere in this circle of
ideas: small positive non-principal eigenvalues and deep negative modes. In the
solved families these are absent or cancel in aggregate (this is precisely the
``collective compensation'' of companion Remark~9.9), which is why the recurrence
runs. The criterion \eqref{eq:teleincr} is strictly weaker to verify than the
full conjecture --- it is a single moment-type inequality with no $g_m$ on the
right --- and it is a clean target for an SOS certificate in the spectral
variables. We record the cleanest sufficient form.

\begin{corollary}[A spectral-gap telescoping criterion]\label{cor:telegap}
If the non-principal spectrum of $T_W$ avoids the open band
$(0,q)\cup(-\tfrac12,-q)$ --- i.e.\ every non-principal eigenvalue lies in
$[-q,0]\cup[q,\tfrac12]$ --- then \eqref{eq:teleincr} holds and the conjecture
holds for $W$ for all odd $m$ (given $D_3\ge0$, which is Goodman's bound,
companion Theorem~3.1).
\end{corollary}

\begin{proof}
Under the hypothesis every summand of \eqref{eq:teleincr} is nonnegative, and
$D_3\ge0$ by the triangle case.
\end{proof}

This is a uniform sufficient condition complementary to
Theorem~\ref{thm:negsafe}: it tolerates eigenvalues in $[q,\tfrac12]$ and in
$[-q,0]$, but forbids the two bands $(0,q)$ and $(-\tfrac12,-q)$. The open
problem is to drop these bands by paying with the principal surplus, i.e.\ to
prove \eqref{eq:teleincr} ``on average over $m$'' rather than termwise --- the
same aggregate target as \eqref{eq:star} and the spectral-excess inequality of
companion \S9.6.

\section{Program P4: Schur convexity and the moment-sequence structure of path
layers}\label{sec:schur}

This sharpens the combinatorial route of companion \S11 by exposing a rigid
structural fact: the path-density sequence is a Hamburger moment sequence, so the
layers are symmetric functions of a totally positive (Hankel) object.

\paragraph{The moment-sequence fact.}
Let $U=1-W$, $x_j=t(P_j,U)=\langle\1,T_U^{\,j}\1\rangle$. Let $\sigma$ be the
scalar spectral measure of the self-adjoint operator $T_U$ at the unit vector
$\1$, supported in $[-1,1]$.

\begin{lemma}[Path densities form a moment sequence]\label{lem:moment}
$x_j=\int_{-1}^{1}t^{\,j}\,d\sigma(t)$ for all $j\ge0$, with
$\int d\sigma=x_0=1$ and $\int t\,d\sigma=x_1=q$. Consequently $(x_j)_{j\ge0}$ is
a Hamburger moment sequence: all Hankel matrices $[x_{i+j}]_{i,j=0}^{N}$ are
positive semidefinite. In particular it is log-convex,
\begin{equation}\label{eq:logconvex}
x_{j-1}\,x_{j+1}\ \ge\ x_j^2\qquad(j\ge1).
\end{equation}
\end{lemma}

\begin{proof}
The spectral theorem applied to $T_U$ and the vector $\1$ gives the integral
representation; $x_0=\|\1\|^2=1$ and $x_1=\langle\1,T_U\1\rangle=q$. Positive
semidefiniteness of the Hankel forms is the defining property of a moment
sequence (for any polynomial $r$, $\sum_{i,j}c_ic_j x_{i+j}=\int |r(t)|^2\,
d\sigma\ge0$ with $r=\sum c_it^i$); \eqref{eq:logconvex} is the $2\times2$ minor,
equivalently Cauchy--Schwarz for $\int t^{j-1}\cdot t^{j+1}d\sigma$.
\end{proof}

\paragraph{Layers as symmetric functions.}
With $n=m-1$, the companion expansion writes
$t(C_m,1-U)-g_m=\sum_{j=0}^{n-1}(-1)^j\tfrac{m}{m-j}E_{n,j}(U)+nE_{n,n}(U)
+(x_n-c_m)$, where the layer is a positive combination of monomials in the
$x_\ell$,
\[
Z_{n,j}(U)=\!\!\sum_{\substack{F\subseteq P_n\\ e(F)=j}}\!\!t(F,U)
=\!\!\sum_{\substack{\text{placements of }j\text{ edges}\\ \text{in }P_n\text{ as disjoint subpaths}}}\!\!\prod_{c}x_{\ell_c},\qquad
E_{n,j}=Z_{n,j}-\binom nj q^{\,j},
\]
the reference value $\binom nj q^j$ being the rank-one specialization
$x_\ell\mapsto q^\ell$. Kim--Lee proved $E_{n,2d}(U)\ge0$ (even layers) by Schur
convexity of $Z_{n,2d}$ in the moment data. The open content is the odd layers
together with the path--cycle gap $x_n-c_m$.

\begin{program}[Master Hankel identity]\label{prog:schur}
Prove a single identity expressing the full alternating sum
$\sum_j(-1)^j\frac{m}{m-j}E_{n,j}+nE_{n,n}+(x_n-c_m)$ as a manifestly
nonnegative combination of (i) Hankel minors of $(x_j)$ (nonnegative by
Lemma~\ref{lem:moment}) and (ii) the deletion gap $x_n-c_m\ge0$. The concrete
lever is to pair each unsigned odd layer $E_{n,2d+1}$ against its even neighbors
using \eqref{eq:logconvex}: log-convexity gives
$x_{2d+1}\le\sqrt{x_{2d}x_{2d+2}}$, which dominates the leading monomial of the
odd layer by a geometric mean of even-layer monomials. The target is a
termwise majorization $|{\rm odd\ layer}|\le\tfrac12({\rm adjacent\ even\
layers})$ making the alternating sum telescope.
\end{program}

\paragraph{Why it could be uniform, and the obstruction.}
The layer recurrence and the moment structure are $m$-independent, so a master
identity is uniform by construction; this is the most ``elementary'' route in the
sense of using only Hankel positivity and one deletion inequality. The risk,
identical to companion \S9.6 in combinatorial dress, is that the odd layers carry
sign information not captured by adjacent even layers --- i.e.\ the pairing is
\emph{not} termwise dominated and only the global sum is controlled, which is the
aggregate phenomenon again. A useful intermediate result, of independent
interest and possibly provable now, is the odd-layer bound
$|E_{n,2d+1}(U)|\le E_{n,2d}(U)^{1/2}E_{n,2d+2}(U)^{1/2}$ as a consequence of
\eqref{eq:logconvex}, which would already control all but the boundary layers.

\section{Program P5: entropy / probabilistic corrected-Sidorenko}\label{sec:entropy}

We set up the route honestly, including the reason it is the longest shot. The
defect to control is $p^m-t(C_m,W)$, and the conjecture asks it to be at most
$p\,q^{m-1}$: a \emph{corrected} Sidorenko bound. (Plain Sidorenko
$t(C_m,W)\ge p^m$ is false for odd $m$, which is the whole difficulty: any
entropy argument must \emph{produce} the negative correction term, not merely
recover $p^m$.)

\paragraph{Homomorphism-measure entropy.}
Let $\nu$ be the probability measure on $[0,1]^{V(C_m)}$ with density
proportional to $\prod_{e=(i,i')}W(x_i,x_{i'})$, normalized by $t(C_m,W)$. The
Gibbs variational principle gives, for any reference product measure $\pi$ with
the correct edge marginals,
\[
\log t(C_m,W)\ \ge\ \mathbb{E}_\pi\!\sum_{e}\log W(x_i,x_{i'})\ -\ \mathrm{KL}(\pi\,\|\,\mathrm{unif}),
\]
and the program is to engineer $\pi$ (a Markov field along the cycle, built from
the transfer operator) so that the right side equals $\log g_m(p)$ up to the
single correction. For \emph{bipartite} $H$ the Conlon--Lee/Szegedy choice of
$\pi$ matches all edge marginals and yields exactly $e(H)\log p$; for odd cycles
the odd girth obstructs a consistent two-colouring of the marginals, leaving an
irreducible mismatch around the cycle --- and that mismatch is precisely the
correction $p\,q^{m-1}$ we must capture rather than discard.

\paragraph{The complement reading and a one-term coupling target.}
The clean structural hint is that $p\,q^{m-1}$ has the exact form ``one edge
present, all $m-1$ complement-edges along the cycle on,'' a \emph{single}
inclusion--exclusion event. This suggests a coupling that produces the defect
from one bad excursion. Equivalently, the slack lost in the companion proofs is
entirely in the crude step $c_m=t(C_m,U)\le x_{m-1}$; an entropy/coupling bound
on the complement cycle density $t(C_m,U)$ directly --- using that $U$ has
density $q<\tfrac12$ and is itself a graphon --- would recover it.

\begin{program}[Entropy control of the complement cycle]\label{prog:entropy}
Produce a probabilistic upper bound $t(C_m,U)\le x_{m-1}-\Delta_m$ with
$\Delta_m\ge0$ large enough that the companion inclusion--exclusion certificate
$\Phi_m$, no longer forced to discard $x_{m-1}-c_m$, becomes nonnegative for all
odd $m$. By companion Proposition~11.1 this retained deletion term is provably
essential, so any successful entropy argument must target $\Delta_m$, not the
leading $p^m$.
\end{program}

\paragraph{Status and a related lever.}
This is the route we are least optimistic about, precisely because odd cycles
fail Sidorenko, so the entropy machine runs against its grain. The one encouraging
fact is that odd cycles are \emph{common} graphs (Sidorenko/folklore): for all
$W$, $t(C_m,W)+t(C_m,1-W)\ge 2^{1-m}$, proved by convexity/Schur arguments akin
to Kim--Lee. Commonality is a different inequality, but it shows the relevant
convexity is available for odd cycles; adapting its proof to the asymmetric,
density-pinned defect $p^m-t(C_m,W)$ is the concrete bridge to attempt.

\section{Program P6: flag algebras / SDP with exact rational rounding}
\label{sec:flag}

A computational route with direct precedent. The constrained minimum
$f_{C_m}(p)=\min\{t(C_m,W):t(K_2,W)=p\}$ is an extremal density problem of exactly
the type flag algebras solve: one writes $t(C_m,\cdot)-g_m(p)$ as a conical
combination of squares of flags plus a nonnegative multiple of the edge-density
constraint, certified by a semidefinite program, and rounds the floating-point
SDP solution to an exact rational certificate in the Pikhurko style. The note's
reference on minimizing pentagons at fixed order and size is precisely this
method applied to $C_5$, and is the natural template.

\begin{program}[Density-uniform flag certificate]\label{prog:flag}
Compute, for each odd $m$ of interest, an exact rational flag-algebra SOS
certificate for $t(C_m,W)\ge g_m(p)$ valid on $p\in[\tfrac12,\tfrac23]$, either by
(i) solving on a fine rational grid in $p$ and proving the certificate
coefficients are sign-definite polynomials in $p$ between grid points (a finite
exact check, compatible with the note's Bernstein/Sturm philosophy), or (ii) a
single ``variable-density'' flag SDP carrying $p$ as a parameter.
\end{program}

\paragraph{Why it can reach where hand certificates cannot, and its ceiling.}
Flag algebras automatically search a far larger space of squares than the
hand-built Gram/Bernstein certificates of the companion note, so they may close
individual cycles (e.g.\ $C_{15},C_{17}$) that defeat the per-cycle hybrid. The
genuine ceiling is \emph{uniformity in $m$}: an SDP is solved at a fixed $m$, so
this route, like the per-cycle hybrid, produces one cycle at a time unless paired
with a structural $m$-induction (Program~P3) or a reduction theorem
(Programs~P1, P2). Its highest value is therefore twofold: as a finisher for
finitely many further cycles, and as an \emph{oracle} --- the exact extremizers it
returns near $p=\tfrac12$ would confirm or refute the near-bipartite structure
that Programs P1--P2 must assume.

\section{Further routes, briefly}\label{sec:further}

For completeness we list routes we judge weaker but not obviously dead.

\begin{enumerate}[leftmargin=2em,label=\textbf{(F\arabic*)}]
\item \textbf{Gradient flow / {\L}ojasiewicz on graphon space.} Run the
$L^2$-gradient flow of $t(C_m,\cdot)$ at fixed $p$; combine the bang--bang
first-order conditions (companion \S10.1) with a {\L}ojasiewicz inequality to
show the flow converges to a low-complexity stationary point, then classify
stationary points. Obstruction: the free Hessian is indefinite (companion
Proposition~10.3), so convergence to the \emph{global} minimum is not automatic;
this is a route to the \emph{structure} of minimizers, feeding P1.
\item \textbf{Tensor-power amplification.} Since $t(C_m,W^{\otimes t})=
t(C_m,W)^t$ while $g_m$ is not multiplicative, a weak per-instance bound does not
tensor up to a strong one directly; the only hope is a self-improving inequality
(if $D_m\ge -\epsilon$ pointwise implies $D_m\ge -\epsilon^2$ via tensoring) which
we have not been able to set up. Listed for completeness; we are pessimistic.
\item \textbf{Interlacing / Cauchy from a rank-one update.} Write
$T_W=T_{W_0}+\text{(rank one)}$ relative to a regular reference $W_0$ of the same
density; Weyl/Cauchy interlacing then bounds how far the spectrum, hence
$\sum\lambda_i^m$, can move. This is a concrete way to attempt
Target~Lemma~\ref{tgt:surplus} (the Perron surplus is the rank-one update's
effect on $\lambda_0$), and is the route we would try first toward
\eqref{eq:star}.
\item \textbf{Real-rootedness / hyperbolic polynomials.} The characteristic
polynomial of the finite truncations of $T_W$ is real-rooted; if the defect
$D_m$ could be written as a coefficient of a hyperbolic polynomial in $(p,q)$ and
the spectral data, then Garding/Br\"and\'en stability theory would give
nonnegativity. Speculative; the difficulty is that $g_m$ mixes $p$ and $q$
nonhomogeneously.
\end{enumerate}

\section{Summary of status}\label{sec:status}

\begin{itemize}[leftmargin=2em]
\item \textbf{Proved, uniform in $m$.}
Theorem~\ref{thm:reform} (resolvent reformulation);
Theorem~\ref{thm:negsafe} (the $T_W\succeq-(1-p)I$ criterion);
Lemma~\ref{lem:onedeep} and Proposition~\ref{prop:onedeep} (one-deep-mode
reduction for $p\ge2/3$);
Lemma~\ref{lem:tele} and Proposition~\ref{prop:tele} with
Corollary~\ref{cor:telegap} (canonical telescoping and the spectral-band
criterion);
Lemma~\ref{lem:moment} (path densities are a Hamburger moment sequence, hence
log-convex). These are independent sufficient conditions and reductions, each
covering a region strictly larger than the regular case.
\item \textbf{Isolated open lemma.} Target Lemma~\ref{tgt:surplus}: a
second-order (no high cycle-power) lower bound on the Perron surplus by the
deep-mode depth. This alone closes $p\ge2/3$ for all odd $m$ via
Proposition~\ref{prop:onedeep}; route (F3) is how we would attack it.
\item \textbf{Six programs, two of which can be uniform in $m$ by their
structure.}
P1 (reflection-positivity symmetrization, \S\ref{sec:rp}) and
P2 (matrix Markov--Krein, \S\ref{sec:markov}) reduce the conjecture to a
low-dimensional family and are uniform by construction;
P3 (telescoping recurrence, \S\ref{sec:telescope}) is uniform via induction and
reduces everything to the single moment inequality \eqref{eq:teleincr};
P4 (Schur convexity / master Hankel identity, \S\ref{sec:schur}) is uniform via
the $m$-independent layer recurrence;
P5 (entropy / corrected Sidorenko, \S\ref{sec:entropy}) is the longest shot,
fighting the failure of Sidorenko for odd cycles;
P6 (flag algebras with exact rounding, \S\ref{sec:flag}) is a per-cycle finisher
and a structural oracle, not uniform on its own.
Routes (F1)--(F4) are listed as weaker but live; (F3) feeds
Target~Lemma~\ref{tgt:surplus} directly.
\item \textbf{Recurring obstruction.} Every clean approach terminates at the same
place: the Perron/non-principal coupling carried by $g$, equivalently the
``aggregate, not modewise'' phenomenon of companion \S9.6. It surfaces as the
lossy Perron split (Remark~\ref{rem:tight}), the unprovable termwise increment
\eqref{eq:teleincr}, the non-termwise odd-vs-even layer pairing
(\S\ref{sec:schur}), and the conditional hypothesis of
Proposition~\ref{prop:onedeep}. We have tried throughout to make this obstruction
quantitative (Target~Lemma~\ref{tgt:surplus}) rather than treat it as folklore.
\item \textbf{Map of dependencies.} A complete uniform proof would most plausibly
combine a \emph{structure} input (P1 or P2, establishing the minimizer is
near-bipartite / few-atom) with a \emph{closing} input on that reduced family
(the low-dimensional inequality of \S\ref{sec:reduction}, or P3/P4 run on the
reduced spectrum). P6 supplies the oracle that tells us which structure to aim
for; (F3) is the most concrete single lemma whose resolution would unlock
$p\ge2/3$.
\end{itemize}

\end{document}